% ============================================================
%  TABLAS DE VALORES DE REFERENCIA — IEC 60335-1
%  Tesis: Equipo de Ensayos de Seguridad Eléctrica Portátil
%  Autor: Valentina Ferrini
% ============================================================
%
% Paquetes requeridos en el preámbulo principal:
%   \usepackage{booktabs}
%   \usepackage{multirow}
%   \usepackage{array}
%   \usepackage{float}
%   \usepackage[table]{xcolor}
%   \usepackage{siunitx}
%   \sisetup{output-decimal-marker={,}}
%   \usepackage[spanish]{babel}
%
% Uso: % ============================================================
%  TABLAS DE VALORES DE REFERENCIA — IEC 60335-1
%  Tesis: Equipo de Ensayos de Seguridad Eléctrica Portátil
%  Autor: Valentina Ferrini
% ============================================================
%
% Paquetes requeridos en el preámbulo principal:
%   \usepackage{booktabs}
%   \usepackage{multirow}
%   \usepackage{array}
%   \usepackage{float}
%   \usepackage[table]{xcolor}
%   \usepackage{siunitx}
%   \sisetup{output-decimal-marker={,}}
%   \usepackage[spanish]{babel}
%
% Uso: % ============================================================
%  TABLAS DE VALORES DE REFERENCIA — IEC 60335-1
%  Tesis: Equipo de Ensayos de Seguridad Eléctrica Portátil
%  Autor: Valentina Ferrini
% ============================================================
%
% Paquetes requeridos en el preámbulo principal:
%   \usepackage{booktabs}
%   \usepackage{multirow}
%   \usepackage{array}
%   \usepackage{float}
%   \usepackage[table]{xcolor}
%   \usepackage{siunitx}
%   \sisetup{output-decimal-marker={,}}
%   \usepackage[spanish]{babel}
%
% Uso: % ============================================================
%  TABLAS DE VALORES DE REFERENCIA — IEC 60335-1
%  Tesis: Equipo de Ensayos de Seguridad Eléctrica Portátil
%  Autor: Valentina Ferrini
% ============================================================
%
% Paquetes requeridos en el preámbulo principal:
%   \usepackage{booktabs}
%   \usepackage{multirow}
%   \usepackage{array}
%   \usepackage{float}
%   \usepackage[table]{xcolor}
%   \usepackage{siunitx}
%   \sisetup{output-decimal-marker={,}}
%   \usepackage[spanish]{babel}
%
% Uso: \input{tablas_IEC60335}
% ============================================================

\section{Valores de referencia según IEC~60335-1}
\label{sec:valores_referencia}

Las tablas presentadas en esta sección sintetizan los valores límite
y los parámetros de ensayo definidos por la norma IEC~60335-1
(\textit{Household and similar electrical appliances -- Safety --
Part~1: General requirements}), edición~4.1 (IEC~60335-1:2001+A1:2004),
para los ensayos contemplados en el equipo desarrollado.
Todos los valores corresponden a aparatos monofásicos con tensión nominal
mayor a \SI{150}{\volt} y hasta \SI{250}{\volt} (\SIrange{220}{240}{\volt},
\SI{50}{\hertz}), rango habitual en el mercado argentino.

La norma distingue dos tipos de ensayo:
\begin{itemize}
    \item \textbf{Ensayos de tipo} (cláusulas~10 a~32): se realizan sobre
          muestras representativas para certificar un diseño.
    \item \textbf{Ensayos de rutina} (Anexo~A, informativo): se realizan
          sobre cada unidad producida para detectar variaciones de
          fabricación que puedan comprometer la seguridad.
\end{itemize}
Por tratarse de un equipo de final de línea, el sistema adopta los
criterios del Anexo~A para rigidez dieléctrica y puesta a tierra.
Para la corriente de fuga y el consumo, que el Anexo~A no contempla,
se adoptan los límites de los ensayos de tipo correspondientes.

% ============================================================
\subsection{Ensayo de consumo}
\label{subsec:tabla_potencia}
% ============================================================

La norma fija la desviación admisible de la potencia
(cláusula~10.1, tabla~1) y de la corriente (cláusula~10.2, tabla~2)
respecto del valor nominal marcado en el aparato.
La tolerancia depende del tipo de aparato y del valor nominal
(tabla~\ref{tab:potencia_desviacion}).
Los aparatos combinados se tratan como aparatos a motor si la potencia
(o corriente) de los motores supera el \SI{50}{\percent} del valor nominal;
en caso contrario, como aparatos de calefacción.

\begin{table}[H]
\centering
\caption{Desviación admisible del consumo respecto del valor nominal --- IEC~60335-1, tablas~1 y~2}
\label{tab:potencia_desviacion}
\renewcommand{\arraystretch}{1.35}
\begin{tabular}{@{} l l l l @{}}
\toprule
\textbf{Tipo de aparato}
    & \textbf{Potencia nominal}
    & \textbf{Corriente nominal}
    & \textbf{Desviación admisible} \\
\midrule
Todos
    & $\leq \SI{25}{\watt}$
    & $\leq \SI{0.2}{\ampere}$
    & $+\SI{20}{\percent}$ \\
\addlinespace
\multirow{2}{*}{Calefacción y combinados}
    & \SIrange[range-phrase={ a }]{25}{200}{\watt}
    & \SIrange[range-phrase={ a }]{0.2}{1}{\ampere}
    & $\pm\SI{10}{\percent}$ \\
    & $> \SI{200}{\watt}$
    & $> \SI{1}{\ampere}$
    & \begin{tabular}[t]{@{}l@{}}
        $+\SI{5}{\percent}$ o $+\SI{20}{\watt}$ ($+\SI{0.10}{\ampere}$),\\
        el mayor; $-\SI{10}{\percent}$
      \end{tabular} \\
\addlinespace
\multirow{2}{*}{A motor}
    & \SIrange[range-phrase={ a }]{25}{300}{\watt}
    & \SIrange[range-phrase={ a }]{0.2}{1.5}{\ampere}
    & $+\SI{20}{\percent}$ \\
    & $> \SI{300}{\watt}$
    & $> \SI{1.5}{\ampere}$
    & \begin{tabular}[t]{@{}l@{}}
        $+\SI{15}{\percent}$ o $+\SI{60}{\watt}$ ($+\SI{0.30}{\ampere}$),\\
        el mayor
      \end{tabular} \\
\bottomrule
\end{tabular}
\end{table}

\begin{table}[H]
\centering
\caption{Condiciones del ensayo de consumo --- IEC~60335-1, cláusulas~5.7, 10.1 y~10.2}
\label{tab:potencia_condiciones}
\renewcommand{\arraystretch}{1.35}
\begin{tabular}{@{} p{4.5cm} p{6cm} p{2.5cm} @{}}
\toprule
\textbf{Parámetro} & \textbf{Condición / valor} & \textbf{Referencia} \\
\midrule
Tensión de alimentación
    & Tensión nominal declarada por el fabricante
    & Cl.~10.1 \\
\addlinespace
Condición de funcionamiento
    & Funcionamiento normal, con todos los circuitos que pueden
      operar simultáneamente en funcionamiento
    & Cl.~10.1 \\
\addlinespace
Momento de la medición
    & Cuando el consumo se ha estabilizado; si varía durante el ciclo,
      se toma el valor medio en un período representativo
    & Cl.~10.1 \\
\addlinespace
Temperatura ambiente
    & \SI{20 \pm 5}{\celsius}, sin corrientes de aire
    & Cl.~5.7 \\
\addlinespace
Magnitud evaluada
    & Potencia activa [\si{\watt}] y corriente eficaz [\si{\ampere}]
    & Cl.~10.1 y~10.2 \\
\bottomrule
\end{tabular}
\end{table}

\noindent Para los aparatos a motor y los de muy bajo consumo la norma
solo fija un límite superior. El sistema agrega un mínimo funcional
\emph{no normativo}, configurable (por defecto \SI{50}{\percent} del valor
nominal), para detectar aparatos que no funcionan.

% ============================================================
\subsection{Ensayo de corriente de fuga}
\label{subsec:tabla_fuga}
% ============================================================

La norma define dos mediciones de corriente de fuga con los mismos
límites (tabla~\ref{tab:fuga_limites}):
\begin{itemize}
    \item \textbf{Fuga a temperatura de funcionamiento} (cláusula~13.2),
          con el aparato funcionando. En este trabajo se la denomina
          \emph{fuga activa}.
    \item \textbf{Fuga con tensión de ensayo} (cláusula~16.2), aplicando
          $\num{1.06}$ veces la tensión nominal sin que el aparato esté en
          funcionamiento normal. En este trabajo se la denomina
          \emph{fuga pasiva}. Solo en este caso los límites se duplican
          bajo ciertas condiciones (tabla~\ref{tab:fuga_duplicacion}).
\end{itemize}

\begin{table}[H]
\centering
\caption{Límites de corriente de fuga --- IEC~60335-1, cláusulas~13.2 y~16.2}
\label{tab:fuga_limites}
\renewcommand{\arraystretch}{1.35}
\begin{tabular}{@{} l r @{}}
\toprule
\textbf{Tipo de aparato} & \textbf{Límite máximo} \\
\midrule
Clase~II
    & \SI{0.25}{\milli\ampere} \\
Clase~0, Clase~0I y Clase~III
    & \SI{0.5}{\milli\ampere} \\
Clase~I portátil
    & \SI{0.75}{\milli\ampere} \\
Clase~I fijo, a motor
    & \SI{3.5}{\milli\ampere} \\
Clase~I fijo, de calefacción
    & \begin{tabular}[t]{@{}r@{}}
        \SI{0.75}{\milli\ampere} o \SI{0.75}{\milli\ampere} por \si{\kilo\watt}\\
        (máximo \SI{5}{\milli\ampere}), el mayor
      \end{tabular} \\
\bottomrule
\end{tabular}
\end{table}

\begin{table}[H]
\centering
\caption{Condiciones que duplican el límite de fuga del ensayo de la cláusula~16.2}
\label{tab:fuga_duplicacion}
\renewcommand{\arraystretch}{1.35}
\begin{tabular}{@{} p{12cm} @{}}
\toprule
\textbf{El límite se duplica si se cumple alguna de las siguientes condiciones} \\
\midrule
Todos los controles tienen posición de apagado en todos los polos. \\
\addlinespace
El aparato no tiene más control que un protector térmico. \\
\addlinespace
Ningún termostato, limitador de temperatura ni regulador de energía
tiene posición de apagado. \\
\addlinespace
El aparato tiene filtros de interferencia de radio; en ese caso, la
corriente de fuga con el filtro desconectado no debe superar el límite
sin duplicar. \\
\bottomrule
\end{tabular}
\end{table}

\begin{table}[H]
\centering
\caption{Condiciones de ensayo para corriente de fuga --- IEC~60335-1, cláusulas~13.1, 13.2 y~16.2}
\label{tab:fuga_condiciones}
\renewcommand{\arraystretch}{1.35}
\begin{tabular}{@{} p{4.5cm} p{8cm} @{}}
\toprule
\textbf{Parámetro de ensayo} & \textbf{Valor / descripción} \\
\midrule
Alimentación (fuga activa, 13.1)
    & Aparatos de calefacción: \num{1.15} veces la potencia nominal. \newline
      Aparatos a motor y combinados: \num{1.06} veces la tensión nominal \\
\addlinespace
Tensión de ensayo (fuga pasiva, 16.2)
    & \num{1.06} veces la tensión nominal (aparatos monofásicos);
      la fuga se mide dentro de los \SI{5}{\second} de aplicada \\
\addlinespace
Red de medición
    & Circuito de la figura~4 de IEC~60990; voltímetro de valor
      eficaz verdadero (\textit{true RMS}) \\
\addlinespace
Puntos de medición
    & Entre cualquier polo de la alimentación y las partes metálicas
      accesibles, incluida una lámina metálica de hasta
      \SI{20 x 10}{\centi\meter} sobre superficies accesibles de
      material aislante \\
\addlinespace
Repetición
    & Se mide con el selector en las posiciones~a y~b
      (ambas polaridades de alimentación) \\
\addlinespace
Aparatos con capacitores e interruptor unipolar
    & Se repite la medición con el interruptor en posición de apagado
      (cláusula~13.2) \\
\bottomrule
\end{tabular}
\end{table}

% ============================================================
\subsection{Ensayo de resistencia de puesta a tierra}
\label{subsec:tabla_tierra}
% ============================================================

\begin{table}[H]
\centering
\caption{Parámetros del ensayo de continuidad a tierra --- IEC~60335-1, cláusula~27.5 y Anexo~A.1}
\label{tab:tierra_referencia}
\renewcommand{\arraystretch}{1.35}
\begin{tabular}{@{} p{3.4cm} p{4.4cm} p{4.4cm} @{}}
\toprule
\textbf{Parámetro}
    & \textbf{Ensayo de tipo (Cl.~27.5)}
    & \textbf{Ensayo de rutina (Anexo~A.1)} \\
\midrule
Corriente de ensayo
    & \num{1.5} veces la corriente nominal o \SI{25}{\ampere}, la mayor
    & Al menos \SI{10}{\ampere}
      (nominal \SI{15}{\ampere} en este equipo) \\
\addlinespace
Tensión de la fuente (en vacío)
    & $\leq \SI{12}{\volt}$ (CA o CC)
    & $\leq \SI{12}{\volt}$ (CA o CC)
      (\SI{2.2}{\volt} CA en este equipo) \\
\addlinespace
Resistencia máxima
    & \SI{100}{\milli\ohm}, sin incluir el cable de alimentación
    & Con cable de alimentación: \SI{200}{\milli\ohm}, o
      \SI{100}{\milli\ohm} más la resistencia del cable. \newline
      Otros aparatos: \SI{100}{\milli\ohm} \\
\addlinespace
Puntos de medición
    & Entre el borne o contacto de tierra y cada parte metálica accesible
    & Entre cada parte metálica accesible puesta a tierra y la
      espiga de tierra de la ficha (o de la ficha del aparato) \\
\addlinespace
Duración
    & En caso de duda, hasta alcanzar condiciones estables
    & Solo el tiempo necesario para medir la caída de tensión \\
\bottomrule
\end{tabular}
\end{table}

\noindent El ensayo solo aplica a aparatos Clase~0I y Clase~I: los aparatos
Clase~0, Clase~II y Clase~III no deben tener provisión para puesta a tierra
(cláusula~27.1).
En ambos casos la norma advierte que la resistencia de contacto entre la
punta de medición y la parte metálica no debe influir en el resultado;
por eso se recomienda el método de cuatro terminales (Kelvin), que elimina
la resistencia de los cables de prueba.

% ============================================================
\subsection{Ensayo de rigidez dieléctrica (tensión resistida)}
\label{subsec:tabla_dielectrica}
% ============================================================

\begin{table}[H]
\centering
\caption{Tensiones de ensayo de rigidez dieléctrica para tensión nominal entre \SI{150}{\volt} y \SI{250}{\volt} --- IEC~60335-1, tablas~4, 7 y~A.1}
\label{tab:hipot_tensiones}
\renewcommand{\arraystretch}{1.35}
\begin{tabular}{@{} l r r r @{}}
\toprule
\textbf{Aislamiento}
    & \textbf{Cl.~13.3}
    & \textbf{Cl.~16.3}
    & \textbf{Anexo~A.2} \\
    & \textbf{(tipo, en caliente)}
    & \textbf{(tipo, tras humedad)}
    & \textbf{(rutina)} \\
\midrule
Básico (Clase~I)
    & \SI{1000}{\volt}
    & \SI{1250}{\volt}
    & \SI{1000}{\volt} \\
Suplementario
    & \SI{1750}{\volt}
    & \SI{1750}{\volt}
    & --- \\
Doble o reforzado (Clase~II)
    & \SI{3000}{\volt}
    & \SI{3000}{\volt}
    & \SI{2500}{\volt} \\
\midrule
Duración
    & \SI{1}{\minute}
    & \SI{1}{\minute}
    & \SI{1}{\second} \\
\bottomrule
\end{tabular}
\end{table}

\begin{table}[H]
\centering
\caption{Condiciones del ensayo de rigidez dieléctrica de rutina --- IEC~60335-1, Anexo~A.2}
\label{tab:hipot_condiciones}
\renewcommand{\arraystretch}{1.35}
\begin{tabular}{@{} p{4.5cm} p{8cm} @{}}
\toprule
\textbf{Parámetro} & \textbf{Valor / descripción} \\
\midrule
Forma de onda
    & Sensiblemente sinusoidal, aproximadamente \SI{50}{\hertz} o
      \SI{60}{\hertz}. Alternativamente, tensión continua de \num{1.5}
      veces el valor de la tabla~\ref{tab:hipot_tensiones} \\
\addlinespace
Duración
    & \SI{1}{\second} \\
\addlinespace
Puntos de aplicación
    & Entre partes activas y partes metálicas accesibles separadas
      de ellas por aislamiento básico (Clase~I) o por aislamiento doble
      o reforzado (Clase~II) \\
\addlinespace
Criterio de aprobación
    & No debe producirse ruptura. Se considera que hay ruptura cuando
      la corriente del circuito de ensayo supera \SI{5}{\milli\ampere};
      el límite puede elevarse hasta \SI{30}{\milli\ampere} para aparatos
      con alta corriente de fuga \\
\addlinespace
Fuente de alta tensión
    & Debe mantener la tensión especificada con la corriente límite y
      contar con un dispositivo que corte al superarla \\
\addlinespace
Condición del aparato
    & Puede ser necesario que el aparato funcione durante el ensayo
      para que la tensión alcance todo el aislamiento (por ejemplo,
      resistencias calefactoras comandadas por relé) \\
\bottomrule
\end{tabular}
\end{table}

\noindent Para los ensayos de tipo (cláusulas~13.3 y~16.3), la fuente de alta
tensión debe poder entregar una corriente de cortocircuito de al menos
\SI{200}{\milli\ampere}, y su protección no debe actuar por debajo de
\SI{100}{\milli\ampere} (tabla~5 de la norma).
El transformador del equipo, con una corriente de cortocircuito limitada
a algunas decenas de \si{\milli\ampere}, es adecuado para el ensayo de
rutina pero no para los ensayos de tipo.

% ============================================================
\subsection{Resumen consolidado de cotas de referencia}
\label{subsec:tabla_resumen}
% ============================================================

La siguiente tabla consolida los límites que aplica el equipo,
facilitando la comparación rápida de los resultados obtenidos frente
a la norma.

\begin{table}[H]
\centering
\caption{Resumen de cotas aplicadas por el equipo --- IEC~60335-1}
\label{tab:resumen_cotas}
\renewcommand{\arraystretch}{1.4}
\rowcolors{2}{gray!10}{white}
\begin{tabular}{@{} l l l l @{}}
\toprule
\textbf{Ensayo}
    & \textbf{Parámetro controlado}
    & \textbf{Límite}
    & \textbf{Referencia} \\
\midrule
Consumo
    & Desviación respecto del valor nominal
    & Según tabla~\ref{tab:potencia_desviacion}
    & Cl.~10.1 y~10.2 \\
Fuga activa
    & $I_{\text{fuga}}$ (Clase~I portátil)
    & $\leq \SI{0.75}{\milli\ampere}$
    & Cl.~13.2 \\
Fuga pasiva
    & $I_{\text{fuga}}$ (Clase~I portátil)
    & \begin{tabular}[t]{@{}l@{}}
        $\leq \SI{0.75}{\milli\ampere}$\\
        ($\leq \SI{1.5}{\milli\ampere}$ con duplicación)
      \end{tabular}
    & Cl.~16.2 \\
Puesta a tierra
    & $R_{\text{PE}}$ con $I \geq \SI{10}{\ampere}$
    & \begin{tabular}[t]{@{}l@{}}
        $\leq \SI{200}{\milli\ohm}$ con cable de alimentación\\
        $\leq \SI{100}{\milli\ohm}$ sin cable
      \end{tabular}
    & Anexo~A.1 \\
Rigidez dieléctrica
    & $I_{\text{HV}}$ a \SI{1000}{\volt} (Clase~I) / \SI{2500}{\volt} (Clase~II)
    & $\leq \SI{5}{\milli\ampere}$ (hasta \SI{30}{\milli\ampere})
    & Anexo~A.2 \\
\bottomrule
\end{tabular}
\end{table}

\begin{table}[H]
\centering
\caption{Parámetros medidos por el equipo en cada ensayo}
\label{tab:parametros_medidos}
\renewcommand{\arraystretch}{1.35}
\begin{tabular}{@{} p{3.5cm} p{8.5cm} @{}}
\toprule
\textbf{Ensayo} & \textbf{Parámetros registrados por el sistema} \\
\midrule
Consumo
    & Tensión $V_{\text{rms}}$ [\si{\volt}], corriente $I_{\text{rms}}$
      [\si{\ampere}], potencia activa $P$ [\si{\watt}], factor de potencia
      $\cos\varphi$, tiempo de ensayo [\si{\second}]; se evalúa el promedio
      del intervalo \\
\addlinespace
Corriente de fuga (pasiva y activa)
    & Corriente de fuga $I_{\text{fuga}}$ [\si{\micro\ampere}],
      tiempo de ensayo [\si{\second}] \\
\addlinespace
Puesta a tierra
    & Corriente de ensayo $I$ [\si{\ampere}] y resistencia calculada
      $R_{\text{PE}}$ [\si{\milli\ohm}]; se evalúa el peor caso del
      intervalo (resistencia máxima y corriente mínima) \\
\addlinespace
Rigidez dieléctrica
    & Tensión aplicada $V_{\text{HV}}$ [\si{\volt}] y corriente en el
      circuito de alta tensión $I_{\text{HV}}$ [\si{\milli\ampere}]; se
      evalúa el peor caso del intervalo (tensión mínima y corriente máxima) \\
\bottomrule
\end{tabular}
\end{table}

\noindent\textbf{Nota sobre la incertidumbre de medición:}
los valores límite de la norma son valores máximos absolutos.
Para garantizar que el equipo no declare como aprobado un DUT
que se encuentra en el margen de tolerancia, se aplica el
criterio de \emph{guardia metrológica}: el sistema emite una
alerta de \emph{valor próximo al límite} cuando el resultado
supera el \SI{80}{\percent} del límite normativo
(p.ej., alerta en $I_{\text{fuga}} > \SI{600}{\micro\ampere}$ para Clase~I),
permitiendo al operario decidir si realiza una medición adicional
con un instrumento patrón antes de emitir el reporte final.

% ============================================================

\section{Valores de referencia según IEC~60335-1}
\label{sec:valores_referencia}

Las tablas presentadas en esta sección sintetizan los valores límite
y los parámetros de ensayo definidos por la norma IEC~60335-1
(\textit{Household and similar electrical appliances -- Safety --
Part~1: General requirements}), edición~4.1 (IEC~60335-1:2001+A1:2004),
para los ensayos contemplados en el equipo desarrollado.
Todos los valores corresponden a aparatos monofásicos con tensión nominal
mayor a \SI{150}{\volt} y hasta \SI{250}{\volt} (\SIrange{220}{240}{\volt},
\SI{50}{\hertz}), rango habitual en el mercado argentino.

La norma distingue dos tipos de ensayo:
\begin{itemize}
    \item \textbf{Ensayos de tipo} (cláusulas~10 a~32): se realizan sobre
          muestras representativas para certificar un diseño.
    \item \textbf{Ensayos de rutina} (Anexo~A, informativo): se realizan
          sobre cada unidad producida para detectar variaciones de
          fabricación que puedan comprometer la seguridad.
\end{itemize}
Por tratarse de un equipo de final de línea, el sistema adopta los
criterios del Anexo~A para rigidez dieléctrica y puesta a tierra.
Para la corriente de fuga y el consumo, que el Anexo~A no contempla,
se adoptan los límites de los ensayos de tipo correspondientes.

% ============================================================
\subsection{Ensayo de consumo}
\label{subsec:tabla_potencia}
% ============================================================

La norma fija la desviación admisible de la potencia
(cláusula~10.1, tabla~1) y de la corriente (cláusula~10.2, tabla~2)
respecto del valor nominal marcado en el aparato.
La tolerancia depende del tipo de aparato y del valor nominal
(tabla~\ref{tab:potencia_desviacion}).
Los aparatos combinados se tratan como aparatos a motor si la potencia
(o corriente) de los motores supera el \SI{50}{\percent} del valor nominal;
en caso contrario, como aparatos de calefacción.

\begin{table}[H]
\centering
\caption{Desviación admisible del consumo respecto del valor nominal --- IEC~60335-1, tablas~1 y~2}
\label{tab:potencia_desviacion}
\renewcommand{\arraystretch}{1.35}
\begin{tabular}{@{} l l l l @{}}
\toprule
\textbf{Tipo de aparato}
    & \textbf{Potencia nominal}
    & \textbf{Corriente nominal}
    & \textbf{Desviación admisible} \\
\midrule
Todos
    & $\leq \SI{25}{\watt}$
    & $\leq \SI{0.2}{\ampere}$
    & $+\SI{20}{\percent}$ \\
\addlinespace
\multirow{2}{*}{Calefacción y combinados}
    & \SIrange[range-phrase={ a }]{25}{200}{\watt}
    & \SIrange[range-phrase={ a }]{0.2}{1}{\ampere}
    & $\pm\SI{10}{\percent}$ \\
    & $> \SI{200}{\watt}$
    & $> \SI{1}{\ampere}$
    & \begin{tabular}[t]{@{}l@{}}
        $+\SI{5}{\percent}$ o $+\SI{20}{\watt}$ ($+\SI{0.10}{\ampere}$),\\
        el mayor; $-\SI{10}{\percent}$
      \end{tabular} \\
\addlinespace
\multirow{2}{*}{A motor}
    & \SIrange[range-phrase={ a }]{25}{300}{\watt}
    & \SIrange[range-phrase={ a }]{0.2}{1.5}{\ampere}
    & $+\SI{20}{\percent}$ \\
    & $> \SI{300}{\watt}$
    & $> \SI{1.5}{\ampere}$
    & \begin{tabular}[t]{@{}l@{}}
        $+\SI{15}{\percent}$ o $+\SI{60}{\watt}$ ($+\SI{0.30}{\ampere}$),\\
        el mayor
      \end{tabular} \\
\bottomrule
\end{tabular}
\end{table}

\begin{table}[H]
\centering
\caption{Condiciones del ensayo de consumo --- IEC~60335-1, cláusulas~5.7, 10.1 y~10.2}
\label{tab:potencia_condiciones}
\renewcommand{\arraystretch}{1.35}
\begin{tabular}{@{} p{4.5cm} p{6cm} p{2.5cm} @{}}
\toprule
\textbf{Parámetro} & \textbf{Condición / valor} & \textbf{Referencia} \\
\midrule
Tensión de alimentación
    & Tensión nominal declarada por el fabricante
    & Cl.~10.1 \\
\addlinespace
Condición de funcionamiento
    & Funcionamiento normal, con todos los circuitos que pueden
      operar simultáneamente en funcionamiento
    & Cl.~10.1 \\
\addlinespace
Momento de la medición
    & Cuando el consumo se ha estabilizado; si varía durante el ciclo,
      se toma el valor medio en un período representativo
    & Cl.~10.1 \\
\addlinespace
Temperatura ambiente
    & \SI{20 \pm 5}{\celsius}, sin corrientes de aire
    & Cl.~5.7 \\
\addlinespace
Magnitud evaluada
    & Potencia activa [\si{\watt}] y corriente eficaz [\si{\ampere}]
    & Cl.~10.1 y~10.2 \\
\bottomrule
\end{tabular}
\end{table}

\noindent Para los aparatos a motor y los de muy bajo consumo la norma
solo fija un límite superior. El sistema agrega un mínimo funcional
\emph{no normativo}, configurable (por defecto \SI{50}{\percent} del valor
nominal), para detectar aparatos que no funcionan.

% ============================================================
\subsection{Ensayo de corriente de fuga}
\label{subsec:tabla_fuga}
% ============================================================

La norma define dos mediciones de corriente de fuga con los mismos
límites (tabla~\ref{tab:fuga_limites}):
\begin{itemize}
    \item \textbf{Fuga a temperatura de funcionamiento} (cláusula~13.2),
          con el aparato funcionando. En este trabajo se la denomina
          \emph{fuga activa}.
    \item \textbf{Fuga con tensión de ensayo} (cláusula~16.2), aplicando
          $\num{1.06}$ veces la tensión nominal sin que el aparato esté en
          funcionamiento normal. En este trabajo se la denomina
          \emph{fuga pasiva}. Solo en este caso los límites se duplican
          bajo ciertas condiciones (tabla~\ref{tab:fuga_duplicacion}).
\end{itemize}

\begin{table}[H]
\centering
\caption{Límites de corriente de fuga --- IEC~60335-1, cláusulas~13.2 y~16.2}
\label{tab:fuga_limites}
\renewcommand{\arraystretch}{1.35}
\begin{tabular}{@{} l r @{}}
\toprule
\textbf{Tipo de aparato} & \textbf{Límite máximo} \\
\midrule
Clase~II
    & \SI{0.25}{\milli\ampere} \\
Clase~0, Clase~0I y Clase~III
    & \SI{0.5}{\milli\ampere} \\
Clase~I portátil
    & \SI{0.75}{\milli\ampere} \\
Clase~I fijo, a motor
    & \SI{3.5}{\milli\ampere} \\
Clase~I fijo, de calefacción
    & \begin{tabular}[t]{@{}r@{}}
        \SI{0.75}{\milli\ampere} o \SI{0.75}{\milli\ampere} por \si{\kilo\watt}\\
        (máximo \SI{5}{\milli\ampere}), el mayor
      \end{tabular} \\
\bottomrule
\end{tabular}
\end{table}

\begin{table}[H]
\centering
\caption{Condiciones que duplican el límite de fuga del ensayo de la cláusula~16.2}
\label{tab:fuga_duplicacion}
\renewcommand{\arraystretch}{1.35}
\begin{tabular}{@{} p{12cm} @{}}
\toprule
\textbf{El límite se duplica si se cumple alguna de las siguientes condiciones} \\
\midrule
Todos los controles tienen posición de apagado en todos los polos. \\
\addlinespace
El aparato no tiene más control que un protector térmico. \\
\addlinespace
Ningún termostato, limitador de temperatura ni regulador de energía
tiene posición de apagado. \\
\addlinespace
El aparato tiene filtros de interferencia de radio; en ese caso, la
corriente de fuga con el filtro desconectado no debe superar el límite
sin duplicar. \\
\bottomrule
\end{tabular}
\end{table}

\begin{table}[H]
\centering
\caption{Condiciones de ensayo para corriente de fuga --- IEC~60335-1, cláusulas~13.1, 13.2 y~16.2}
\label{tab:fuga_condiciones}
\renewcommand{\arraystretch}{1.35}
\begin{tabular}{@{} p{4.5cm} p{8cm} @{}}
\toprule
\textbf{Parámetro de ensayo} & \textbf{Valor / descripción} \\
\midrule
Alimentación (fuga activa, 13.1)
    & Aparatos de calefacción: \num{1.15} veces la potencia nominal. \newline
      Aparatos a motor y combinados: \num{1.06} veces la tensión nominal \\
\addlinespace
Tensión de ensayo (fuga pasiva, 16.2)
    & \num{1.06} veces la tensión nominal (aparatos monofásicos);
      la fuga se mide dentro de los \SI{5}{\second} de aplicada \\
\addlinespace
Red de medición
    & Circuito de la figura~4 de IEC~60990; voltímetro de valor
      eficaz verdadero (\textit{true RMS}) \\
\addlinespace
Puntos de medición
    & Entre cualquier polo de la alimentación y las partes metálicas
      accesibles, incluida una lámina metálica de hasta
      \SI{20 x 10}{\centi\meter} sobre superficies accesibles de
      material aislante \\
\addlinespace
Repetición
    & Se mide con el selector en las posiciones~a y~b
      (ambas polaridades de alimentación) \\
\addlinespace
Aparatos con capacitores e interruptor unipolar
    & Se repite la medición con el interruptor en posición de apagado
      (cláusula~13.2) \\
\bottomrule
\end{tabular}
\end{table}

% ============================================================
\subsection{Ensayo de resistencia de puesta a tierra}
\label{subsec:tabla_tierra}
% ============================================================

\begin{table}[H]
\centering
\caption{Parámetros del ensayo de continuidad a tierra --- IEC~60335-1, cláusula~27.5 y Anexo~A.1}
\label{tab:tierra_referencia}
\renewcommand{\arraystretch}{1.35}
\begin{tabular}{@{} p{3.4cm} p{4.4cm} p{4.4cm} @{}}
\toprule
\textbf{Parámetro}
    & \textbf{Ensayo de tipo (Cl.~27.5)}
    & \textbf{Ensayo de rutina (Anexo~A.1)} \\
\midrule
Corriente de ensayo
    & \num{1.5} veces la corriente nominal o \SI{25}{\ampere}, la mayor
    & Al menos \SI{10}{\ampere}
      (nominal \SI{15}{\ampere} en este equipo) \\
\addlinespace
Tensión de la fuente (en vacío)
    & $\leq \SI{12}{\volt}$ (CA o CC)
    & $\leq \SI{12}{\volt}$ (CA o CC)
      (\SI{2.2}{\volt} CA en este equipo) \\
\addlinespace
Resistencia máxima
    & \SI{100}{\milli\ohm}, sin incluir el cable de alimentación
    & Con cable de alimentación: \SI{200}{\milli\ohm}, o
      \SI{100}{\milli\ohm} más la resistencia del cable. \newline
      Otros aparatos: \SI{100}{\milli\ohm} \\
\addlinespace
Puntos de medición
    & Entre el borne o contacto de tierra y cada parte metálica accesible
    & Entre cada parte metálica accesible puesta a tierra y la
      espiga de tierra de la ficha (o de la ficha del aparato) \\
\addlinespace
Duración
    & En caso de duda, hasta alcanzar condiciones estables
    & Solo el tiempo necesario para medir la caída de tensión \\
\bottomrule
\end{tabular}
\end{table}

\noindent El ensayo solo aplica a aparatos Clase~0I y Clase~I: los aparatos
Clase~0, Clase~II y Clase~III no deben tener provisión para puesta a tierra
(cláusula~27.1).
En ambos casos la norma advierte que la resistencia de contacto entre la
punta de medición y la parte metálica no debe influir en el resultado;
por eso se recomienda el método de cuatro terminales (Kelvin), que elimina
la resistencia de los cables de prueba.

% ============================================================
\subsection{Ensayo de rigidez dieléctrica (tensión resistida)}
\label{subsec:tabla_dielectrica}
% ============================================================

\begin{table}[H]
\centering
\caption{Tensiones de ensayo de rigidez dieléctrica para tensión nominal entre \SI{150}{\volt} y \SI{250}{\volt} --- IEC~60335-1, tablas~4, 7 y~A.1}
\label{tab:hipot_tensiones}
\renewcommand{\arraystretch}{1.35}
\begin{tabular}{@{} l r r r @{}}
\toprule
\textbf{Aislamiento}
    & \textbf{Cl.~13.3}
    & \textbf{Cl.~16.3}
    & \textbf{Anexo~A.2} \\
    & \textbf{(tipo, en caliente)}
    & \textbf{(tipo, tras humedad)}
    & \textbf{(rutina)} \\
\midrule
Básico (Clase~I)
    & \SI{1000}{\volt}
    & \SI{1250}{\volt}
    & \SI{1000}{\volt} \\
Suplementario
    & \SI{1750}{\volt}
    & \SI{1750}{\volt}
    & --- \\
Doble o reforzado (Clase~II)
    & \SI{3000}{\volt}
    & \SI{3000}{\volt}
    & \SI{2500}{\volt} \\
\midrule
Duración
    & \SI{1}{\minute}
    & \SI{1}{\minute}
    & \SI{1}{\second} \\
\bottomrule
\end{tabular}
\end{table}

\begin{table}[H]
\centering
\caption{Condiciones del ensayo de rigidez dieléctrica de rutina --- IEC~60335-1, Anexo~A.2}
\label{tab:hipot_condiciones}
\renewcommand{\arraystretch}{1.35}
\begin{tabular}{@{} p{4.5cm} p{8cm} @{}}
\toprule
\textbf{Parámetro} & \textbf{Valor / descripción} \\
\midrule
Forma de onda
    & Sensiblemente sinusoidal, aproximadamente \SI{50}{\hertz} o
      \SI{60}{\hertz}. Alternativamente, tensión continua de \num{1.5}
      veces el valor de la tabla~\ref{tab:hipot_tensiones} \\
\addlinespace
Duración
    & \SI{1}{\second} \\
\addlinespace
Puntos de aplicación
    & Entre partes activas y partes metálicas accesibles separadas
      de ellas por aislamiento básico (Clase~I) o por aislamiento doble
      o reforzado (Clase~II) \\
\addlinespace
Criterio de aprobación
    & No debe producirse ruptura. Se considera que hay ruptura cuando
      la corriente del circuito de ensayo supera \SI{5}{\milli\ampere};
      el límite puede elevarse hasta \SI{30}{\milli\ampere} para aparatos
      con alta corriente de fuga \\
\addlinespace
Fuente de alta tensión
    & Debe mantener la tensión especificada con la corriente límite y
      contar con un dispositivo que corte al superarla \\
\addlinespace
Condición del aparato
    & Puede ser necesario que el aparato funcione durante el ensayo
      para que la tensión alcance todo el aislamiento (por ejemplo,
      resistencias calefactoras comandadas por relé) \\
\bottomrule
\end{tabular}
\end{table}

\noindent Para los ensayos de tipo (cláusulas~13.3 y~16.3), la fuente de alta
tensión debe poder entregar una corriente de cortocircuito de al menos
\SI{200}{\milli\ampere}, y su protección no debe actuar por debajo de
\SI{100}{\milli\ampere} (tabla~5 de la norma).
El transformador del equipo, con una corriente de cortocircuito limitada
a algunas decenas de \si{\milli\ampere}, es adecuado para el ensayo de
rutina pero no para los ensayos de tipo.

% ============================================================
\subsection{Resumen consolidado de cotas de referencia}
\label{subsec:tabla_resumen}
% ============================================================

La siguiente tabla consolida los límites que aplica el equipo,
facilitando la comparación rápida de los resultados obtenidos frente
a la norma.

\begin{table}[H]
\centering
\caption{Resumen de cotas aplicadas por el equipo --- IEC~60335-1}
\label{tab:resumen_cotas}
\renewcommand{\arraystretch}{1.4}
\rowcolors{2}{gray!10}{white}
\begin{tabular}{@{} l l l l @{}}
\toprule
\textbf{Ensayo}
    & \textbf{Parámetro controlado}
    & \textbf{Límite}
    & \textbf{Referencia} \\
\midrule
Consumo
    & Desviación respecto del valor nominal
    & Según tabla~\ref{tab:potencia_desviacion}
    & Cl.~10.1 y~10.2 \\
Fuga activa
    & $I_{\text{fuga}}$ (Clase~I portátil)
    & $\leq \SI{0.75}{\milli\ampere}$
    & Cl.~13.2 \\
Fuga pasiva
    & $I_{\text{fuga}}$ (Clase~I portátil)
    & \begin{tabular}[t]{@{}l@{}}
        $\leq \SI{0.75}{\milli\ampere}$\\
        ($\leq \SI{1.5}{\milli\ampere}$ con duplicación)
      \end{tabular}
    & Cl.~16.2 \\
Puesta a tierra
    & $R_{\text{PE}}$ con $I \geq \SI{10}{\ampere}$
    & \begin{tabular}[t]{@{}l@{}}
        $\leq \SI{200}{\milli\ohm}$ con cable de alimentación\\
        $\leq \SI{100}{\milli\ohm}$ sin cable
      \end{tabular}
    & Anexo~A.1 \\
Rigidez dieléctrica
    & $I_{\text{HV}}$ a \SI{1000}{\volt} (Clase~I) / \SI{2500}{\volt} (Clase~II)
    & $\leq \SI{5}{\milli\ampere}$ (hasta \SI{30}{\milli\ampere})
    & Anexo~A.2 \\
\bottomrule
\end{tabular}
\end{table}

\begin{table}[H]
\centering
\caption{Parámetros medidos por el equipo en cada ensayo}
\label{tab:parametros_medidos}
\renewcommand{\arraystretch}{1.35}
\begin{tabular}{@{} p{3.5cm} p{8.5cm} @{}}
\toprule
\textbf{Ensayo} & \textbf{Parámetros registrados por el sistema} \\
\midrule
Consumo
    & Tensión $V_{\text{rms}}$ [\si{\volt}], corriente $I_{\text{rms}}$
      [\si{\ampere}], potencia activa $P$ [\si{\watt}], factor de potencia
      $\cos\varphi$, tiempo de ensayo [\si{\second}]; se evalúa el promedio
      del intervalo \\
\addlinespace
Corriente de fuga (pasiva y activa)
    & Corriente de fuga $I_{\text{fuga}}$ [\si{\micro\ampere}],
      tiempo de ensayo [\si{\second}] \\
\addlinespace
Puesta a tierra
    & Corriente de ensayo $I$ [\si{\ampere}] y resistencia calculada
      $R_{\text{PE}}$ [\si{\milli\ohm}]; se evalúa el peor caso del
      intervalo (resistencia máxima y corriente mínima) \\
\addlinespace
Rigidez dieléctrica
    & Tensión aplicada $V_{\text{HV}}$ [\si{\volt}] y corriente en el
      circuito de alta tensión $I_{\text{HV}}$ [\si{\milli\ampere}]; se
      evalúa el peor caso del intervalo (tensión mínima y corriente máxima) \\
\bottomrule
\end{tabular}
\end{table}

\noindent\textbf{Nota sobre la incertidumbre de medición:}
los valores límite de la norma son valores máximos absolutos.
Para garantizar que el equipo no declare como aprobado un DUT
que se encuentra en el margen de tolerancia, se aplica el
criterio de \emph{guardia metrológica}: el sistema emite una
alerta de \emph{valor próximo al límite} cuando el resultado
supera el \SI{80}{\percent} del límite normativo
(p.ej., alerta en $I_{\text{fuga}} > \SI{600}{\micro\ampere}$ para Clase~I),
permitiendo al operario decidir si realiza una medición adicional
con un instrumento patrón antes de emitir el reporte final.

% ============================================================

\section{Valores de referencia según IEC~60335-1}
\label{sec:valores_referencia}

Las tablas presentadas en esta sección sintetizan los valores límite
y los parámetros de ensayo definidos por la norma IEC~60335-1
(\textit{Household and similar electrical appliances -- Safety --
Part~1: General requirements}), edición~4.1 (IEC~60335-1:2001+A1:2004),
para los ensayos contemplados en el equipo desarrollado.
Todos los valores corresponden a aparatos monofásicos con tensión nominal
mayor a \SI{150}{\volt} y hasta \SI{250}{\volt} (\SIrange{220}{240}{\volt},
\SI{50}{\hertz}), rango habitual en el mercado argentino.

La norma distingue dos tipos de ensayo:
\begin{itemize}
    \item \textbf{Ensayos de tipo} (cláusulas~10 a~32): se realizan sobre
          muestras representativas para certificar un diseño.
    \item \textbf{Ensayos de rutina} (Anexo~A, informativo): se realizan
          sobre cada unidad producida para detectar variaciones de
          fabricación que puedan comprometer la seguridad.
\end{itemize}
Por tratarse de un equipo de final de línea, el sistema adopta los
criterios del Anexo~A para rigidez dieléctrica y puesta a tierra.
Para la corriente de fuga y el consumo, que el Anexo~A no contempla,
se adoptan los límites de los ensayos de tipo correspondientes.

% ============================================================
\subsection{Ensayo de consumo}
\label{subsec:tabla_potencia}
% ============================================================

La norma fija la desviación admisible de la potencia
(cláusula~10.1, tabla~1) y de la corriente (cláusula~10.2, tabla~2)
respecto del valor nominal marcado en el aparato.
La tolerancia depende del tipo de aparato y del valor nominal
(tabla~\ref{tab:potencia_desviacion}).
Los aparatos combinados se tratan como aparatos a motor si la potencia
(o corriente) de los motores supera el \SI{50}{\percent} del valor nominal;
en caso contrario, como aparatos de calefacción.

\begin{table}[H]
\centering
\caption{Desviación admisible del consumo respecto del valor nominal --- IEC~60335-1, tablas~1 y~2}
\label{tab:potencia_desviacion}
\renewcommand{\arraystretch}{1.35}
\begin{tabular}{@{} l l l l @{}}
\toprule
\textbf{Tipo de aparato}
    & \textbf{Potencia nominal}
    & \textbf{Corriente nominal}
    & \textbf{Desviación admisible} \\
\midrule
Todos
    & $\leq \SI{25}{\watt}$
    & $\leq \SI{0.2}{\ampere}$
    & $+\SI{20}{\percent}$ \\
\addlinespace
\multirow{2}{*}{Calefacción y combinados}
    & \SIrange[range-phrase={ a }]{25}{200}{\watt}
    & \SIrange[range-phrase={ a }]{0.2}{1}{\ampere}
    & $\pm\SI{10}{\percent}$ \\
    & $> \SI{200}{\watt}$
    & $> \SI{1}{\ampere}$
    & \begin{tabular}[t]{@{}l@{}}
        $+\SI{5}{\percent}$ o $+\SI{20}{\watt}$ ($+\SI{0.10}{\ampere}$),\\
        el mayor; $-\SI{10}{\percent}$
      \end{tabular} \\
\addlinespace
\multirow{2}{*}{A motor}
    & \SIrange[range-phrase={ a }]{25}{300}{\watt}
    & \SIrange[range-phrase={ a }]{0.2}{1.5}{\ampere}
    & $+\SI{20}{\percent}$ \\
    & $> \SI{300}{\watt}$
    & $> \SI{1.5}{\ampere}$
    & \begin{tabular}[t]{@{}l@{}}
        $+\SI{15}{\percent}$ o $+\SI{60}{\watt}$ ($+\SI{0.30}{\ampere}$),\\
        el mayor
      \end{tabular} \\
\bottomrule
\end{tabular}
\end{table}

\begin{table}[H]
\centering
\caption{Condiciones del ensayo de consumo --- IEC~60335-1, cláusulas~5.7, 10.1 y~10.2}
\label{tab:potencia_condiciones}
\renewcommand{\arraystretch}{1.35}
\begin{tabular}{@{} p{4.5cm} p{6cm} p{2.5cm} @{}}
\toprule
\textbf{Parámetro} & \textbf{Condición / valor} & \textbf{Referencia} \\
\midrule
Tensión de alimentación
    & Tensión nominal declarada por el fabricante
    & Cl.~10.1 \\
\addlinespace
Condición de funcionamiento
    & Funcionamiento normal, con todos los circuitos que pueden
      operar simultáneamente en funcionamiento
    & Cl.~10.1 \\
\addlinespace
Momento de la medición
    & Cuando el consumo se ha estabilizado; si varía durante el ciclo,
      se toma el valor medio en un período representativo
    & Cl.~10.1 \\
\addlinespace
Temperatura ambiente
    & \SI{20 \pm 5}{\celsius}, sin corrientes de aire
    & Cl.~5.7 \\
\addlinespace
Magnitud evaluada
    & Potencia activa [\si{\watt}] y corriente eficaz [\si{\ampere}]
    & Cl.~10.1 y~10.2 \\
\bottomrule
\end{tabular}
\end{table}

\noindent Para los aparatos a motor y los de muy bajo consumo la norma
solo fija un límite superior. El sistema agrega un mínimo funcional
\emph{no normativo}, configurable (por defecto \SI{50}{\percent} del valor
nominal), para detectar aparatos que no funcionan.

% ============================================================
\subsection{Ensayo de corriente de fuga}
\label{subsec:tabla_fuga}
% ============================================================

La norma define dos mediciones de corriente de fuga con los mismos
límites (tabla~\ref{tab:fuga_limites}):
\begin{itemize}
    \item \textbf{Fuga a temperatura de funcionamiento} (cláusula~13.2),
          con el aparato funcionando. En este trabajo se la denomina
          \emph{fuga activa}.
    \item \textbf{Fuga con tensión de ensayo} (cláusula~16.2), aplicando
          $\num{1.06}$ veces la tensión nominal sin que el aparato esté en
          funcionamiento normal. En este trabajo se la denomina
          \emph{fuga pasiva}. Solo en este caso los límites se duplican
          bajo ciertas condiciones (tabla~\ref{tab:fuga_duplicacion}).
\end{itemize}

\begin{table}[H]
\centering
\caption{Límites de corriente de fuga --- IEC~60335-1, cláusulas~13.2 y~16.2}
\label{tab:fuga_limites}
\renewcommand{\arraystretch}{1.35}
\begin{tabular}{@{} l r @{}}
\toprule
\textbf{Tipo de aparato} & \textbf{Límite máximo} \\
\midrule
Clase~II
    & \SI{0.25}{\milli\ampere} \\
Clase~0, Clase~0I y Clase~III
    & \SI{0.5}{\milli\ampere} \\
Clase~I portátil
    & \SI{0.75}{\milli\ampere} \\
Clase~I fijo, a motor
    & \SI{3.5}{\milli\ampere} \\
Clase~I fijo, de calefacción
    & \begin{tabular}[t]{@{}r@{}}
        \SI{0.75}{\milli\ampere} o \SI{0.75}{\milli\ampere} por \si{\kilo\watt}\\
        (máximo \SI{5}{\milli\ampere}), el mayor
      \end{tabular} \\
\bottomrule
\end{tabular}
\end{table}

\begin{table}[H]
\centering
\caption{Condiciones que duplican el límite de fuga del ensayo de la cláusula~16.2}
\label{tab:fuga_duplicacion}
\renewcommand{\arraystretch}{1.35}
\begin{tabular}{@{} p{12cm} @{}}
\toprule
\textbf{El límite se duplica si se cumple alguna de las siguientes condiciones} \\
\midrule
Todos los controles tienen posición de apagado en todos los polos. \\
\addlinespace
El aparato no tiene más control que un protector térmico. \\
\addlinespace
Ningún termostato, limitador de temperatura ni regulador de energía
tiene posición de apagado. \\
\addlinespace
El aparato tiene filtros de interferencia de radio; en ese caso, la
corriente de fuga con el filtro desconectado no debe superar el límite
sin duplicar. \\
\bottomrule
\end{tabular}
\end{table}

\begin{table}[H]
\centering
\caption{Condiciones de ensayo para corriente de fuga --- IEC~60335-1, cláusulas~13.1, 13.2 y~16.2}
\label{tab:fuga_condiciones}
\renewcommand{\arraystretch}{1.35}
\begin{tabular}{@{} p{4.5cm} p{8cm} @{}}
\toprule
\textbf{Parámetro de ensayo} & \textbf{Valor / descripción} \\
\midrule
Alimentación (fuga activa, 13.1)
    & Aparatos de calefacción: \num{1.15} veces la potencia nominal. \newline
      Aparatos a motor y combinados: \num{1.06} veces la tensión nominal \\
\addlinespace
Tensión de ensayo (fuga pasiva, 16.2)
    & \num{1.06} veces la tensión nominal (aparatos monofásicos);
      la fuga se mide dentro de los \SI{5}{\second} de aplicada \\
\addlinespace
Red de medición
    & Circuito de la figura~4 de IEC~60990; voltímetro de valor
      eficaz verdadero (\textit{true RMS}) \\
\addlinespace
Puntos de medición
    & Entre cualquier polo de la alimentación y las partes metálicas
      accesibles, incluida una lámina metálica de hasta
      \SI{20 x 10}{\centi\meter} sobre superficies accesibles de
      material aislante \\
\addlinespace
Repetición
    & Se mide con el selector en las posiciones~a y~b
      (ambas polaridades de alimentación) \\
\addlinespace
Aparatos con capacitores e interruptor unipolar
    & Se repite la medición con el interruptor en posición de apagado
      (cláusula~13.2) \\
\bottomrule
\end{tabular}
\end{table}

% ============================================================
\subsection{Ensayo de resistencia de puesta a tierra}
\label{subsec:tabla_tierra}
% ============================================================

\begin{table}[H]
\centering
\caption{Parámetros del ensayo de continuidad a tierra --- IEC~60335-1, cláusula~27.5 y Anexo~A.1}
\label{tab:tierra_referencia}
\renewcommand{\arraystretch}{1.35}
\begin{tabular}{@{} p{3.4cm} p{4.4cm} p{4.4cm} @{}}
\toprule
\textbf{Parámetro}
    & \textbf{Ensayo de tipo (Cl.~27.5)}
    & \textbf{Ensayo de rutina (Anexo~A.1)} \\
\midrule
Corriente de ensayo
    & \num{1.5} veces la corriente nominal o \SI{25}{\ampere}, la mayor
    & Al menos \SI{10}{\ampere}
      (nominal \SI{15}{\ampere} en este equipo) \\
\addlinespace
Tensión de la fuente (en vacío)
    & $\leq \SI{12}{\volt}$ (CA o CC)
    & $\leq \SI{12}{\volt}$ (CA o CC)
      (\SI{2.2}{\volt} CA en este equipo) \\
\addlinespace
Resistencia máxima
    & \SI{100}{\milli\ohm}, sin incluir el cable de alimentación
    & Con cable de alimentación: \SI{200}{\milli\ohm}, o
      \SI{100}{\milli\ohm} más la resistencia del cable. \newline
      Otros aparatos: \SI{100}{\milli\ohm} \\
\addlinespace
Puntos de medición
    & Entre el borne o contacto de tierra y cada parte metálica accesible
    & Entre cada parte metálica accesible puesta a tierra y la
      espiga de tierra de la ficha (o de la ficha del aparato) \\
\addlinespace
Duración
    & En caso de duda, hasta alcanzar condiciones estables
    & Solo el tiempo necesario para medir la caída de tensión \\
\bottomrule
\end{tabular}
\end{table}

\noindent El ensayo solo aplica a aparatos Clase~0I y Clase~I: los aparatos
Clase~0, Clase~II y Clase~III no deben tener provisión para puesta a tierra
(cláusula~27.1).
En ambos casos la norma advierte que la resistencia de contacto entre la
punta de medición y la parte metálica no debe influir en el resultado;
por eso se recomienda el método de cuatro terminales (Kelvin), que elimina
la resistencia de los cables de prueba.

% ============================================================
\subsection{Ensayo de rigidez dieléctrica (tensión resistida)}
\label{subsec:tabla_dielectrica}
% ============================================================

\begin{table}[H]
\centering
\caption{Tensiones de ensayo de rigidez dieléctrica para tensión nominal entre \SI{150}{\volt} y \SI{250}{\volt} --- IEC~60335-1, tablas~4, 7 y~A.1}
\label{tab:hipot_tensiones}
\renewcommand{\arraystretch}{1.35}
\begin{tabular}{@{} l r r r @{}}
\toprule
\textbf{Aislamiento}
    & \textbf{Cl.~13.3}
    & \textbf{Cl.~16.3}
    & \textbf{Anexo~A.2} \\
    & \textbf{(tipo, en caliente)}
    & \textbf{(tipo, tras humedad)}
    & \textbf{(rutina)} \\
\midrule
Básico (Clase~I)
    & \SI{1000}{\volt}
    & \SI{1250}{\volt}
    & \SI{1000}{\volt} \\
Suplementario
    & \SI{1750}{\volt}
    & \SI{1750}{\volt}
    & --- \\
Doble o reforzado (Clase~II)
    & \SI{3000}{\volt}
    & \SI{3000}{\volt}
    & \SI{2500}{\volt} \\
\midrule
Duración
    & \SI{1}{\minute}
    & \SI{1}{\minute}
    & \SI{1}{\second} \\
\bottomrule
\end{tabular}
\end{table}

\begin{table}[H]
\centering
\caption{Condiciones del ensayo de rigidez dieléctrica de rutina --- IEC~60335-1, Anexo~A.2}
\label{tab:hipot_condiciones}
\renewcommand{\arraystretch}{1.35}
\begin{tabular}{@{} p{4.5cm} p{8cm} @{}}
\toprule
\textbf{Parámetro} & \textbf{Valor / descripción} \\
\midrule
Forma de onda
    & Sensiblemente sinusoidal, aproximadamente \SI{50}{\hertz} o
      \SI{60}{\hertz}. Alternativamente, tensión continua de \num{1.5}
      veces el valor de la tabla~\ref{tab:hipot_tensiones} \\
\addlinespace
Duración
    & \SI{1}{\second} \\
\addlinespace
Puntos de aplicación
    & Entre partes activas y partes metálicas accesibles separadas
      de ellas por aislamiento básico (Clase~I) o por aislamiento doble
      o reforzado (Clase~II) \\
\addlinespace
Criterio de aprobación
    & No debe producirse ruptura. Se considera que hay ruptura cuando
      la corriente del circuito de ensayo supera \SI{5}{\milli\ampere};
      el límite puede elevarse hasta \SI{30}{\milli\ampere} para aparatos
      con alta corriente de fuga \\
\addlinespace
Fuente de alta tensión
    & Debe mantener la tensión especificada con la corriente límite y
      contar con un dispositivo que corte al superarla \\
\addlinespace
Condición del aparato
    & Puede ser necesario que el aparato funcione durante el ensayo
      para que la tensión alcance todo el aislamiento (por ejemplo,
      resistencias calefactoras comandadas por relé) \\
\bottomrule
\end{tabular}
\end{table}

\noindent Para los ensayos de tipo (cláusulas~13.3 y~16.3), la fuente de alta
tensión debe poder entregar una corriente de cortocircuito de al menos
\SI{200}{\milli\ampere}, y su protección no debe actuar por debajo de
\SI{100}{\milli\ampere} (tabla~5 de la norma).
El transformador del equipo, con una corriente de cortocircuito limitada
a algunas decenas de \si{\milli\ampere}, es adecuado para el ensayo de
rutina pero no para los ensayos de tipo.

% ============================================================
\subsection{Resumen consolidado de cotas de referencia}
\label{subsec:tabla_resumen}
% ============================================================

La siguiente tabla consolida los límites que aplica el equipo,
facilitando la comparación rápida de los resultados obtenidos frente
a la norma.

\begin{table}[H]
\centering
\caption{Resumen de cotas aplicadas por el equipo --- IEC~60335-1}
\label{tab:resumen_cotas}
\renewcommand{\arraystretch}{1.4}
\rowcolors{2}{gray!10}{white}
\begin{tabular}{@{} l l l l @{}}
\toprule
\textbf{Ensayo}
    & \textbf{Parámetro controlado}
    & \textbf{Límite}
    & \textbf{Referencia} \\
\midrule
Consumo
    & Desviación respecto del valor nominal
    & Según tabla~\ref{tab:potencia_desviacion}
    & Cl.~10.1 y~10.2 \\
Fuga activa
    & $I_{\text{fuga}}$ (Clase~I portátil)
    & $\leq \SI{0.75}{\milli\ampere}$
    & Cl.~13.2 \\
Fuga pasiva
    & $I_{\text{fuga}}$ (Clase~I portátil)
    & \begin{tabular}[t]{@{}l@{}}
        $\leq \SI{0.75}{\milli\ampere}$\\
        ($\leq \SI{1.5}{\milli\ampere}$ con duplicación)
      \end{tabular}
    & Cl.~16.2 \\
Puesta a tierra
    & $R_{\text{PE}}$ con $I \geq \SI{10}{\ampere}$
    & \begin{tabular}[t]{@{}l@{}}
        $\leq \SI{200}{\milli\ohm}$ con cable de alimentación\\
        $\leq \SI{100}{\milli\ohm}$ sin cable
      \end{tabular}
    & Anexo~A.1 \\
Rigidez dieléctrica
    & $I_{\text{HV}}$ a \SI{1000}{\volt} (Clase~I) / \SI{2500}{\volt} (Clase~II)
    & $\leq \SI{5}{\milli\ampere}$ (hasta \SI{30}{\milli\ampere})
    & Anexo~A.2 \\
\bottomrule
\end{tabular}
\end{table}

\begin{table}[H]
\centering
\caption{Parámetros medidos por el equipo en cada ensayo}
\label{tab:parametros_medidos}
\renewcommand{\arraystretch}{1.35}
\begin{tabular}{@{} p{3.5cm} p{8.5cm} @{}}
\toprule
\textbf{Ensayo} & \textbf{Parámetros registrados por el sistema} \\
\midrule
Consumo
    & Tensión $V_{\text{rms}}$ [\si{\volt}], corriente $I_{\text{rms}}$
      [\si{\ampere}], potencia activa $P$ [\si{\watt}], factor de potencia
      $\cos\varphi$, tiempo de ensayo [\si{\second}]; se evalúa el promedio
      del intervalo \\
\addlinespace
Corriente de fuga (pasiva y activa)
    & Corriente de fuga $I_{\text{fuga}}$ [\si{\micro\ampere}],
      tiempo de ensayo [\si{\second}] \\
\addlinespace
Puesta a tierra
    & Corriente de ensayo $I$ [\si{\ampere}] y resistencia calculada
      $R_{\text{PE}}$ [\si{\milli\ohm}]; se evalúa el peor caso del
      intervalo (resistencia máxima y corriente mínima) \\
\addlinespace
Rigidez dieléctrica
    & Tensión aplicada $V_{\text{HV}}$ [\si{\volt}] y corriente en el
      circuito de alta tensión $I_{\text{HV}}$ [\si{\milli\ampere}]; se
      evalúa el peor caso del intervalo (tensión mínima y corriente máxima) \\
\bottomrule
\end{tabular}
\end{table}

\noindent\textbf{Nota sobre la incertidumbre de medición:}
los valores límite de la norma son valores máximos absolutos.
Para garantizar que el equipo no declare como aprobado un DUT
que se encuentra en el margen de tolerancia, se aplica el
criterio de \emph{guardia metrológica}: el sistema emite una
alerta de \emph{valor próximo al límite} cuando el resultado
supera el \SI{80}{\percent} del límite normativo
(p.ej., alerta en $I_{\text{fuga}} > \SI{600}{\micro\ampere}$ para Clase~I),
permitiendo al operario decidir si realiza una medición adicional
con un instrumento patrón antes de emitir el reporte final.

% ============================================================

\section{Valores de referencia según IEC~60335-1}
\label{sec:valores_referencia}

Las tablas presentadas en esta sección sintetizan los valores límite
y los parámetros de ensayo definidos por la norma IEC~60335-1
(\textit{Household and similar electrical appliances -- Safety --
Part~1: General requirements}), edición~4.1 (IEC~60335-1:2001+A1:2004),
para los ensayos contemplados en el equipo desarrollado.
Todos los valores corresponden a aparatos monofásicos con tensión nominal
mayor a \SI{150}{\volt} y hasta \SI{250}{\volt} (\SIrange{220}{240}{\volt},
\SI{50}{\hertz}), rango habitual en el mercado argentino.

La norma distingue dos tipos de ensayo:
\begin{itemize}
    \item \textbf{Ensayos de tipo} (cláusulas~10 a~32): se realizan sobre
          muestras representativas para certificar un diseño.
    \item \textbf{Ensayos de rutina} (Anexo~A, informativo): se realizan
          sobre cada unidad producida para detectar variaciones de
          fabricación que puedan comprometer la seguridad.
\end{itemize}
Por tratarse de un equipo de final de línea, el sistema adopta los
criterios del Anexo~A para rigidez dieléctrica y puesta a tierra.
Para la corriente de fuga y el consumo, que el Anexo~A no contempla,
se adoptan los límites de los ensayos de tipo correspondientes.

% ============================================================
\subsection{Ensayo de consumo}
\label{subsec:tabla_potencia}
% ============================================================

La norma fija la desviación admisible de la potencia
(cláusula~10.1, tabla~1) y de la corriente (cláusula~10.2, tabla~2)
respecto del valor nominal marcado en el aparato.
La tolerancia depende del tipo de aparato y del valor nominal
(tabla~\ref{tab:potencia_desviacion}).
Los aparatos combinados se tratan como aparatos a motor si la potencia
(o corriente) de los motores supera el \SI{50}{\percent} del valor nominal;
en caso contrario, como aparatos de calefacción.

\begin{table}[H]
\centering
\caption{Desviación admisible del consumo respecto del valor nominal --- IEC~60335-1, tablas~1 y~2}
\label{tab:potencia_desviacion}
\renewcommand{\arraystretch}{1.35}
\begin{tabular}{@{} l l l l @{}}
\toprule
\textbf{Tipo de aparato}
    & \textbf{Potencia nominal}
    & \textbf{Corriente nominal}
    & \textbf{Desviación admisible} \\
\midrule
Todos
    & $\leq \SI{25}{\watt}$
    & $\leq \SI{0.2}{\ampere}$
    & $+\SI{20}{\percent}$ \\
\addlinespace
\multirow{2}{*}{Calefacción y combinados}
    & \SIrange[range-phrase={ a }]{25}{200}{\watt}
    & \SIrange[range-phrase={ a }]{0.2}{1}{\ampere}
    & $\pm\SI{10}{\percent}$ \\
    & $> \SI{200}{\watt}$
    & $> \SI{1}{\ampere}$
    & \begin{tabular}[t]{@{}l@{}}
        $+\SI{5}{\percent}$ o $+\SI{20}{\watt}$ ($+\SI{0.10}{\ampere}$),\\
        el mayor; $-\SI{10}{\percent}$
      \end{tabular} \\
\addlinespace
\multirow{2}{*}{A motor}
    & \SIrange[range-phrase={ a }]{25}{300}{\watt}
    & \SIrange[range-phrase={ a }]{0.2}{1.5}{\ampere}
    & $+\SI{20}{\percent}$ \\
    & $> \SI{300}{\watt}$
    & $> \SI{1.5}{\ampere}$
    & \begin{tabular}[t]{@{}l@{}}
        $+\SI{15}{\percent}$ o $+\SI{60}{\watt}$ ($+\SI{0.30}{\ampere}$),\\
        el mayor
      \end{tabular} \\
\bottomrule
\end{tabular}
\end{table}

\begin{table}[H]
\centering
\caption{Condiciones del ensayo de consumo --- IEC~60335-1, cláusulas~5.7, 10.1 y~10.2}
\label{tab:potencia_condiciones}
\renewcommand{\arraystretch}{1.35}
\begin{tabular}{@{} p{4.5cm} p{6cm} p{2.5cm} @{}}
\toprule
\textbf{Parámetro} & \textbf{Condición / valor} & \textbf{Referencia} \\
\midrule
Tensión de alimentación
    & Tensión nominal declarada por el fabricante
    & Cl.~10.1 \\
\addlinespace
Condición de funcionamiento
    & Funcionamiento normal, con todos los circuitos que pueden
      operar simultáneamente en funcionamiento
    & Cl.~10.1 \\
\addlinespace
Momento de la medición
    & Cuando el consumo se ha estabilizado; si varía durante el ciclo,
      se toma el valor medio en un período representativo
    & Cl.~10.1 \\
\addlinespace
Temperatura ambiente
    & \SI{20 \pm 5}{\celsius}, sin corrientes de aire
    & Cl.~5.7 \\
\addlinespace
Magnitud evaluada
    & Potencia activa [\si{\watt}] y corriente eficaz [\si{\ampere}]
    & Cl.~10.1 y~10.2 \\
\bottomrule
\end{tabular}
\end{table}

\noindent Para los aparatos a motor y los de muy bajo consumo la norma
solo fija un límite superior. El sistema agrega un mínimo funcional
\emph{no normativo}, configurable (por defecto \SI{50}{\percent} del valor
nominal), para detectar aparatos que no funcionan.

% ============================================================
\subsection{Ensayo de corriente de fuga}
\label{subsec:tabla_fuga}
% ============================================================

La norma define dos mediciones de corriente de fuga con los mismos
límites (tabla~\ref{tab:fuga_limites}):
\begin{itemize}
    \item \textbf{Fuga a temperatura de funcionamiento} (cláusula~13.2),
          con el aparato funcionando. En este trabajo se la denomina
          \emph{fuga activa}.
    \item \textbf{Fuga con tensión de ensayo} (cláusula~16.2), aplicando
          $\num{1.06}$ veces la tensión nominal sin que el aparato esté en
          funcionamiento normal. En este trabajo se la denomina
          \emph{fuga pasiva}. Solo en este caso los límites se duplican
          bajo ciertas condiciones (tabla~\ref{tab:fuga_duplicacion}).
\end{itemize}

\begin{table}[H]
\centering
\caption{Límites de corriente de fuga --- IEC~60335-1, cláusulas~13.2 y~16.2}
\label{tab:fuga_limites}
\renewcommand{\arraystretch}{1.35}
\begin{tabular}{@{} l r @{}}
\toprule
\textbf{Tipo de aparato} & \textbf{Límite máximo} \\
\midrule
Clase~II
    & \SI{0.25}{\milli\ampere} \\
Clase~0, Clase~0I y Clase~III
    & \SI{0.5}{\milli\ampere} \\
Clase~I portátil
    & \SI{0.75}{\milli\ampere} \\
Clase~I fijo, a motor
    & \SI{3.5}{\milli\ampere} \\
Clase~I fijo, de calefacción
    & \begin{tabular}[t]{@{}r@{}}
        \SI{0.75}{\milli\ampere} o \SI{0.75}{\milli\ampere} por \si{\kilo\watt}\\
        (máximo \SI{5}{\milli\ampere}), el mayor
      \end{tabular} \\
\bottomrule
\end{tabular}
\end{table}

\begin{table}[H]
\centering
\caption{Condiciones que duplican el límite de fuga del ensayo de la cláusula~16.2}
\label{tab:fuga_duplicacion}
\renewcommand{\arraystretch}{1.35}
\begin{tabular}{@{} p{12cm} @{}}
\toprule
\textbf{El límite se duplica si se cumple alguna de las siguientes condiciones} \\
\midrule
Todos los controles tienen posición de apagado en todos los polos. \\
\addlinespace
El aparato no tiene más control que un protector térmico. \\
\addlinespace
Ningún termostato, limitador de temperatura ni regulador de energía
tiene posición de apagado. \\
\addlinespace
El aparato tiene filtros de interferencia de radio; en ese caso, la
corriente de fuga con el filtro desconectado no debe superar el límite
sin duplicar. \\
\bottomrule
\end{tabular}
\end{table}

\begin{table}[H]
\centering
\caption{Condiciones de ensayo para corriente de fuga --- IEC~60335-1, cláusulas~13.1, 13.2 y~16.2}
\label{tab:fuga_condiciones}
\renewcommand{\arraystretch}{1.35}
\begin{tabular}{@{} p{4.5cm} p{8cm} @{}}
\toprule
\textbf{Parámetro de ensayo} & \textbf{Valor / descripción} \\
\midrule
Alimentación (fuga activa, 13.1)
    & Aparatos de calefacción: \num{1.15} veces la potencia nominal. \newline
      Aparatos a motor y combinados: \num{1.06} veces la tensión nominal \\
\addlinespace
Tensión de ensayo (fuga pasiva, 16.2)
    & \num{1.06} veces la tensión nominal (aparatos monofásicos);
      la fuga se mide dentro de los \SI{5}{\second} de aplicada \\
\addlinespace
Red de medición
    & Circuito de la figura~4 de IEC~60990; voltímetro de valor
      eficaz verdadero (\textit{true RMS}) \\
\addlinespace
Puntos de medición
    & Entre cualquier polo de la alimentación y las partes metálicas
      accesibles, incluida una lámina metálica de hasta
      \SI{20 x 10}{\centi\meter} sobre superficies accesibles de
      material aislante \\
\addlinespace
Repetición
    & Se mide con el selector en las posiciones~a y~b
      (ambas polaridades de alimentación) \\
\addlinespace
Aparatos con capacitores e interruptor unipolar
    & Se repite la medición con el interruptor en posición de apagado
      (cláusula~13.2) \\
\bottomrule
\end{tabular}
\end{table}

% ============================================================
\subsection{Ensayo de resistencia de puesta a tierra}
\label{subsec:tabla_tierra}
% ============================================================

\begin{table}[H]
\centering
\caption{Parámetros del ensayo de continuidad a tierra --- IEC~60335-1, cláusula~27.5 y Anexo~A.1}
\label{tab:tierra_referencia}
\renewcommand{\arraystretch}{1.35}
\begin{tabular}{@{} p{3.4cm} p{4.4cm} p{4.4cm} @{}}
\toprule
\textbf{Parámetro}
    & \textbf{Ensayo de tipo (Cl.~27.5)}
    & \textbf{Ensayo de rutina (Anexo~A.1)} \\
\midrule
Corriente de ensayo
    & \num{1.5} veces la corriente nominal o \SI{25}{\ampere}, la mayor
    & Al menos \SI{10}{\ampere}
      (nominal \SI{15}{\ampere} en este equipo) \\
\addlinespace
Tensión de la fuente (en vacío)
    & $\leq \SI{12}{\volt}$ (CA o CC)
    & $\leq \SI{12}{\volt}$ (CA o CC)
      (\SI{2.2}{\volt} CA en este equipo) \\
\addlinespace
Resistencia máxima
    & \SI{100}{\milli\ohm}, sin incluir el cable de alimentación
    & Con cable de alimentación: \SI{200}{\milli\ohm}, o
      \SI{100}{\milli\ohm} más la resistencia del cable. \newline
      Otros aparatos: \SI{100}{\milli\ohm} \\
\addlinespace
Puntos de medición
    & Entre el borne o contacto de tierra y cada parte metálica accesible
    & Entre cada parte metálica accesible puesta a tierra y la
      espiga de tierra de la ficha (o de la ficha del aparato) \\
\addlinespace
Duración
    & En caso de duda, hasta alcanzar condiciones estables
    & Solo el tiempo necesario para medir la caída de tensión \\
\bottomrule
\end{tabular}
\end{table}

\noindent El ensayo solo aplica a aparatos Clase~0I y Clase~I: los aparatos
Clase~0, Clase~II y Clase~III no deben tener provisión para puesta a tierra
(cláusula~27.1).
En ambos casos la norma advierte que la resistencia de contacto entre la
punta de medición y la parte metálica no debe influir en el resultado;
por eso se recomienda el método de cuatro terminales (Kelvin), que elimina
la resistencia de los cables de prueba.

% ============================================================
\subsection{Ensayo de rigidez dieléctrica (tensión resistida)}
\label{subsec:tabla_dielectrica}
% ============================================================

\begin{table}[H]
\centering
\caption{Tensiones de ensayo de rigidez dieléctrica para tensión nominal entre \SI{150}{\volt} y \SI{250}{\volt} --- IEC~60335-1, tablas~4, 7 y~A.1}
\label{tab:hipot_tensiones}
\renewcommand{\arraystretch}{1.35}
\begin{tabular}{@{} l r r r @{}}
\toprule
\textbf{Aislamiento}
    & \textbf{Cl.~13.3}
    & \textbf{Cl.~16.3}
    & \textbf{Anexo~A.2} \\
    & \textbf{(tipo, en caliente)}
    & \textbf{(tipo, tras humedad)}
    & \textbf{(rutina)} \\
\midrule
Básico (Clase~I)
    & \SI{1000}{\volt}
    & \SI{1250}{\volt}
    & \SI{1000}{\volt} \\
Suplementario
    & \SI{1750}{\volt}
    & \SI{1750}{\volt}
    & --- \\
Doble o reforzado (Clase~II)
    & \SI{3000}{\volt}
    & \SI{3000}{\volt}
    & \SI{2500}{\volt} \\
\midrule
Duración
    & \SI{1}{\minute}
    & \SI{1}{\minute}
    & \SI{1}{\second} \\
\bottomrule
\end{tabular}
\end{table}

\begin{table}[H]
\centering
\caption{Condiciones del ensayo de rigidez dieléctrica de rutina --- IEC~60335-1, Anexo~A.2}
\label{tab:hipot_condiciones}
\renewcommand{\arraystretch}{1.35}
\begin{tabular}{@{} p{4.5cm} p{8cm} @{}}
\toprule
\textbf{Parámetro} & \textbf{Valor / descripción} \\
\midrule
Forma de onda
    & Sensiblemente sinusoidal, aproximadamente \SI{50}{\hertz} o
      \SI{60}{\hertz}. Alternativamente, tensión continua de \num{1.5}
      veces el valor de la tabla~\ref{tab:hipot_tensiones} \\
\addlinespace
Duración
    & \SI{1}{\second} \\
\addlinespace
Puntos de aplicación
    & Entre partes activas y partes metálicas accesibles separadas
      de ellas por aislamiento básico (Clase~I) o por aislamiento doble
      o reforzado (Clase~II) \\
\addlinespace
Criterio de aprobación
    & No debe producirse ruptura. Se considera que hay ruptura cuando
      la corriente del circuito de ensayo supera \SI{5}{\milli\ampere};
      el límite puede elevarse hasta \SI{30}{\milli\ampere} para aparatos
      con alta corriente de fuga \\
\addlinespace
Fuente de alta tensión
    & Debe mantener la tensión especificada con la corriente límite y
      contar con un dispositivo que corte al superarla \\
\addlinespace
Condición del aparato
    & Puede ser necesario que el aparato funcione durante el ensayo
      para que la tensión alcance todo el aislamiento (por ejemplo,
      resistencias calefactoras comandadas por relé) \\
\bottomrule
\end{tabular}
\end{table}

\noindent Para los ensayos de tipo (cláusulas~13.3 y~16.3), la fuente de alta
tensión debe poder entregar una corriente de cortocircuito de al menos
\SI{200}{\milli\ampere}, y su protección no debe actuar por debajo de
\SI{100}{\milli\ampere} (tabla~5 de la norma).
El transformador del equipo, con una corriente de cortocircuito limitada
a algunas decenas de \si{\milli\ampere}, es adecuado para el ensayo de
rutina pero no para los ensayos de tipo.

% ============================================================
\subsection{Resumen consolidado de cotas de referencia}
\label{subsec:tabla_resumen}
% ============================================================

La siguiente tabla consolida los límites que aplica el equipo,
facilitando la comparación rápida de los resultados obtenidos frente
a la norma.

\begin{table}[H]
\centering
\caption{Resumen de cotas aplicadas por el equipo --- IEC~60335-1}
\label{tab:resumen_cotas}
\renewcommand{\arraystretch}{1.4}
\rowcolors{2}{gray!10}{white}
\begin{tabular}{@{} l l l l @{}}
\toprule
\textbf{Ensayo}
    & \textbf{Parámetro controlado}
    & \textbf{Límite}
    & \textbf{Referencia} \\
\midrule
Consumo
    & Desviación respecto del valor nominal
    & Según tabla~\ref{tab:potencia_desviacion}
    & Cl.~10.1 y~10.2 \\
Fuga activa
    & $I_{\text{fuga}}$ (Clase~I portátil)
    & $\leq \SI{0.75}{\milli\ampere}$
    & Cl.~13.2 \\
Fuga pasiva
    & $I_{\text{fuga}}$ (Clase~I portátil)
    & \begin{tabular}[t]{@{}l@{}}
        $\leq \SI{0.75}{\milli\ampere}$\\
        ($\leq \SI{1.5}{\milli\ampere}$ con duplicación)
      \end{tabular}
    & Cl.~16.2 \\
Puesta a tierra
    & $R_{\text{PE}}$ con $I \geq \SI{10}{\ampere}$
    & \begin{tabular}[t]{@{}l@{}}
        $\leq \SI{200}{\milli\ohm}$ con cable de alimentación\\
        $\leq \SI{100}{\milli\ohm}$ sin cable
      \end{tabular}
    & Anexo~A.1 \\
Rigidez dieléctrica
    & $I_{\text{HV}}$ a \SI{1000}{\volt} (Clase~I) / \SI{2500}{\volt} (Clase~II)
    & $\leq \SI{5}{\milli\ampere}$ (hasta \SI{30}{\milli\ampere})
    & Anexo~A.2 \\
\bottomrule
\end{tabular}
\end{table}

\begin{table}[H]
\centering
\caption{Parámetros medidos por el equipo en cada ensayo}
\label{tab:parametros_medidos}
\renewcommand{\arraystretch}{1.35}
\begin{tabular}{@{} p{3.5cm} p{8.5cm} @{}}
\toprule
\textbf{Ensayo} & \textbf{Parámetros registrados por el sistema} \\
\midrule
Consumo
    & Tensión $V_{\text{rms}}$ [\si{\volt}], corriente $I_{\text{rms}}$
      [\si{\ampere}], potencia activa $P$ [\si{\watt}], factor de potencia
      $\cos\varphi$, tiempo de ensayo [\si{\second}]; se evalúa el promedio
      del intervalo \\
\addlinespace
Corriente de fuga (pasiva y activa)
    & Corriente de fuga $I_{\text{fuga}}$ [\si{\micro\ampere}],
      tiempo de ensayo [\si{\second}] \\
\addlinespace
Puesta a tierra
    & Corriente de ensayo $I$ [\si{\ampere}] y resistencia calculada
      $R_{\text{PE}}$ [\si{\milli\ohm}]; se evalúa el peor caso del
      intervalo (resistencia máxima y corriente mínima) \\
\addlinespace
Rigidez dieléctrica
    & Tensión aplicada $V_{\text{HV}}$ [\si{\volt}] y corriente en el
      circuito de alta tensión $I_{\text{HV}}$ [\si{\milli\ampere}]; se
      evalúa el peor caso del intervalo (tensión mínima y corriente máxima) \\
\bottomrule
\end{tabular}
\end{table}

\noindent\textbf{Nota sobre la incertidumbre de medición:}
los valores límite de la norma son valores máximos absolutos.
Para garantizar que el equipo no declare como aprobado un DUT
que se encuentra en el margen de tolerancia, se aplica el
criterio de \emph{guardia metrológica}: el sistema emite una
alerta de \emph{valor próximo al límite} cuando el resultado
supera el \SI{80}{\percent} del límite normativo
(p.ej., alerta en $I_{\text{fuga}} > \SI{600}{\micro\ampere}$ para Clase~I),
permitiendo al operario decidir si realiza una medición adicional
con un instrumento patrón antes de emitir el reporte final.
